\chapter{Introduction}
\label{ch:introduction}

\section{Context and motivations}\label{sec:context-and-motivations}
Design pattern detection has become a crucial aspect of software engineering, as it supports quality assurance,
technical debt management, code smell identification, and overall maintainability.
This process is closely related to refactoring, a systematic method to updating code without changing its
external behavior, enhancing its readability and structure.

Code smells typically arise from codebase growth or poor programming practices, acting as indicators of
underlying issues that can be addressed during refactoring: while they are often easy to find and fix, they may be just
symptoms of deeper archithectural flaws.
Technical debt refers to the maintenance cost growth due to suboptimal design or implementation choices: although in the
short term they can speed up code development, over time they make the system more complex to understand, modify, test,
and evolve.

Design patterns can be used as a means to deal with structural issues and guide refactoring decisions, as they represent
a solution to common problems in software development. Since patterns are described at a high level of abstraction, there
are several ways to choose and apply them, leading to different concrete implementations.
While useful, design patterns should not be applied indiscriminately, since forcing their application in unsuitable
contexts may lead to overengineering, as well as an increase in technical debt.

However, the identification and application of design patterns are not fully deterministic processes, since they rely on
the project context, the considered abstraction level, and the developer interpretation.
As a result, it may not be possible to establish a universally valid answer to the question of which design pattern to
use: such an answer is inevitably influenced by subjective factors.

In particular, LLMs may help to solve this problem, since they can help automate and speed up refactoring, despite
the subjective nature of the task. Design pattern detection requires, in fact, a semantic and contextual understanding
of the source code: this need makes the use of such models particularly relevant, precisely because they do not merely
search for rigid patterns, but are able to analyze structure, class responsibilities, and the relationships between
components.
This process represents an application of artificial intelligence to a branch of software engineering, named
\underline{AI4SE}.

LLMs are promising tools, but their ability to support design pattern recommendation must be evaluated systematically
in order to understand their reliability in realistic scenarios: a software engineer should be able to assess whether
a model is suggesting an appropriate pattern. As mentioned before, the impossibility of establishing a fully universal
answer makes the correctness of a recommended design pattern difficult to define; therefore, suitable performance
indicators must be identified.


\section{Objectives and contributions}\label{sec:objectives-and-contributions}
The main goal of this thesis consists of exploring how effective LLM-based design pattern detection is today in realistic
scenarios. In order to achieve this, the identification of a way to evaluate LLMs' performance is crucial because of the
need to understand software architecture and suggest the best design pattern. This analysis will investigate local model
configurations by identifying, for each chosen LLM, the best parameter settings and subjecting them to realistic
scenarios. The applicability of a specific pattern depends on multiple factors, and it is not guaranteed: the models
will be challenged by architectural ambiguity.

The main purpose in terms of implementation is to create an end-to-end tool, which will support refactoring by taking
as input Python code and providing a refactored version that applies the most suitable design pattern.
The work involves an extension of the \underline{LLM Design Pattern Refactoring}
project~\cite{zisa2024llmdesignpatternrefactoring}, which enables the application of a previously selected design
pattern by implementing a component that suggests the most suitable design pattern given the code.

This thesis makes the following contributions:
\begin{itemize}
    \item Adding new code examples for design patterns that were not present in the dataset of the previous project. Each example has been provided with a refactored version, implementing a specific design pattern, and a baseline version, where the design pattern has been removed. Baseline and refactored versions are functionally equivalent.
    \item Labeling of each baseline example in the dataset with a list of applicable design patterns, ranked by applicability.
    \item Implementation of a tool for querying a local LLM, aimed at generating a list of design patterns ranked according to their applicability to the analyzed code. Each query is characterized by the specification of the LLM name, textual prompt, and model temperature.
    \item Evaluation of LLM predictions according to seven metrics designed to evaluate the ranking quality of model predictions.
    \item Conduct of an ablation study aimed at establishing which query yields the best results.
\end{itemize}

\section{Thesis Organization}\label{sec:thesis-organization}
The thesis is structured as follows:
\begin{itemize}
    \item \emph{Chapter 2} provides the necessary concepts to understand design pattern definitions and how to interact with Large Language Models.
    \item \emph{Chapter 3} describes the dataset, explaining the refactored examples and which design pattern they implement.
    \item \emph{Chapter 4} introduces the LLM-based design pattern detection tool.
    \item \emph{Chapter 5} defines the evaluation methods for model predictions and discusses the results in light of them.
    \item \emph{Chapter 6} discusses related work.
    \item \emph{Chapter 7} concludes the thesis by summarizing the main results, discussing the possibilities and limitations of the approach.
\end{itemize}




